\title{Large Language Models Can Automatically Engineer Features for Few-Shot Tabular Learning}

\begin{abstract}
Large Language Models (LLMs) have demonstrated a significant aptitude for handling complex, novel reasoning tasks, making them especially promising for tabular learning, which is fundamental to a broad array of real-world use cases. This work introduces a new in-context learning paradigm, \textsf{FeatLLM}, which harnesses LLMs as automated feature constructors to transform the original input dataset into representations that are more amenable to tabular prediction tasks. The engineered features facilitate class probability estimation through straightforward downstream machine learning algorithms, such as linear regression, thereby supporting efficient and highly effective few-shot learning. At inference, \textsf{FeatLLM} relies solely on this simple predictive model in conjunction with the constructed features, omitting the need for LLM interaction with each test instance. In contrast to prior LLM-driven strategies, \textsf{FeatLLM} dispenses with recurrent LLM queries for individual predictions and requires only API access to the LLM, circumventing issues like prompt length constraints. Our extensive evaluation over various tabular datasets from multiple domains demonstrates that \textsf{FeatLLM} produces superior features and high-quality decision rules, offering substantial improvements (averaging 10\%) over established approaches such as TabLLM and STUNT.
\end{abstract}

\section{Introduction}
Large Language Models (LLMs) have recently exhibited extraordinary proficiency in generating coherent responses to tasks they have not encountered during training and do so without necessitating additional fine-tuning for each specific task~\cite{creswell2022selection,imani2023mathprompter,taylor2022galactica}. The vast prior knowledge encoded during pretraining on large-scale text datasets empowers LLMs to serve as a proxy for human expertise across a range of domains~\cite{Genomes2021,singhal2023large,wilcox2003role}, thereby facilitating automation in fields spanning from dialogue comprehension~\cite{finch2023leveraging} to automated program repair~\cite{fan2023automated}. By embedding concise task descriptions and a few exemplar data points in their prompts, LLMs can discern relevant contextual information and focus selectively on salient features and cases~\cite{brown2020language,zhao2023survey}. This behavior underlines the possibility of leveraging LLMs as proficient feature engineers.

The exploitation of LLMs’ inherent domain knowledge and logical inference abilities has recently been expanded to tabular learning problems. For example, leveraging information like the correlation between age and increased disease risk can substantially aid in disease classification. Contemporary techniques typically express the tasks and variables using natural language, serialize the structured data, and present it to an LLM~\cite{dinh2022lift,wang2023anypredict}. This input preparation is commonly followed by adopting in-context learning strategies~\cite{nam2023semi} or applying parameter-efficient fine-tuning methods~\cite{dinh2022lift,hegselmann2023tabllm}. Harnessing LLMs’ domain prior has proven advantageous, frequently surpassing conventional tabular learning benchmarks, especially in settings with limited labeled data~\cite{hegselmann2023tabllm}.

Existing LLM-based approaches for tabular learning exhibit several key drawbacks. Notably, end-to-end prediction systems require executing at least one LLM forward pass per data instance, resulting in significant computational overhead. Furthermore, achieving high predictive performance often necessitates fine-tuning, which is impractical for LLMs lacking publicly available training interfaces or fine-tuning APIs—a pressing issue since recent state-of-the-art LLMs generally restrict usage to inference APIs only~\cite{anil2023palm,openai2023gpt4}. In addition, such methods are often ill-suited to processing longer prompts~\cite{liu2023lost}; as the number of tabular features grows and the text representation increases in length, both feasibility and performance may deteriorate. These issues fundamentally impede the deployment of LLM-driven tabular learning techniques in practical environments.

To address these challenges, our method departs from the traditional paradigm of using LLMs exclusively for direct prediction. Instead, we employ LLMs to conduct feature engineering, taking advantage of their pretrained knowledge in a more strategic fashion. Specifically, we introduce an in-context learning paradigm, \textsf{FeatLLM}, designed to uncover the “criteria” or decision rules used by LLMs during classification or prediction tasks. For example, when detecting a particular medical condition, the LLM can articulate rule-based associations specifying which feature values indicate the presence of the disease. These extracted criteria inform the creation of novel features to replace or supplement the original set in subsequent predictive modeling. By shifting the LLM’s primary function from prediction to feature extraction, we enable the use of simpler models downstream; once features have been distilled, sophisticated machine learning algorithms are no longer necessary at inference time, thereby enhancing efficiency. The in-context learning capabilities of LLMs allow us to generate these features directly via prompt engineering—simply by supplying well-chosen exemplars—without further parameter updates. Additionally, strategies inspired by ensemble techniques~\cite{breiman2001random}, such as feature bagging, can be utilized to circumvent issues arising from unwieldy prompt lengths.

\textsf{FeatLLM} decomposes the prompt formulation process for automated feature engineering into two primary components: (1) identifying and interpreting the problem at hand, alongside establishing connections between feature sets and target variables; and (2) utilizing insights acquired from this analysis, together with limited in-context examples, to develop critical decision rules for class discrimination. Through this structured mode of reasoning, the large language models are guided to ignore irrelevant or redundant features during the rule formulation phase. Subsequently, these extracted rules are implemented on the data (interpreted as code instructions), leading to the transformation of raw data into binary indicators reflecting rule adherence per instance. To predict class membership probabilities based on these derived features, a simple and interpretable model is trained. To further enhance model reliability, an ensemble methodology is adopted, with iterative feature creation combined with feature bagging techniques. Proper calibration of the number of selected features and sample size during bagging helps alleviate challenges posed by overly long prompts. Importantly, \textsf{FeatLLM} does not require explicit model training and maintains minimal inference overhead, making its deployment with LLMs practical for diverse applied contexts.

An extensive empirical study of \textsf{FeatLLM} is conducted across 13 tabular datasets under conditions of limited data availability, demonstrating consistent and high performance. Our experiments confirm that large language models are adept at leveraging both their embedded prior knowledge and insights gleaned directly from the dataset, resulting in the construction of effective features. When compared with existing few-shot learning approaches, our methodology consistently yields superior results. All source code for replication is made accessible through an anonymized GitHub repository at~\texttt{https://github.com/Sungwon-Han/FeatLLM}.

\section{Related Work}

\subsection{Few-Shot Learning with Tabular Data}  
Few-shot learning has significantly advanced the ability to derive robust, generalizable feature representations from data, despite only having limited labeled instances~\cite{chen2018closer}. While early research in this area concentrated on visual data, few-shot paradigms have since demonstrated adaptability to multiple data types, such as text, audio, and more recently, tabular formats~\cite{majumder2022few,nam2023stunt,schick2021exploiting}. However, utilizing sparse labeled data in learning introduces susceptibility to capturing spurious dependencies~\cite{bartlett2021deep}. To mitigate this, recent methodologies have integrated supplementary information sources or injected domain-specific prior knowledge, which serves to instill inductive biases favorable for model training. Notable strategies include leveraging abundant unlabeled data sets with strategic data augmentations in semi-supervised learning scenarios~\cite{ucar2021subtab,yoon2020vime} and employing task-agnostic, unsupervised pre-training for model initialization~\cite{somepalli2021saint}. These unsupervised tasks frequently require partially masking or corrupting input data, subsequently reconstructing the original information as the learning objective~\cite{arik2021tabnet,majmundar2022met}, or adopting data augmentation techniques within contrastive learning frameworks~\cite{bahri2022scarf,chen2020simple}. Nevertheless, due to the diversity inherent in tabular data, it remains challenging to craft augmentation policies that are broadly effective, and such techniques often fail to deliver substantial improvements in few-shot conditions~\cite{nam2023stunt}.

A line of recent studies has suggested expanding the training data pool to encompass a variety of tabular datasets—even those unrelated to the eventual downstream application—and applying transfer learning to adapt models to new tabular tasks~\cite{zhang2023generative,levin2022transfer,zhu2023xtab}. As an example, TransTab~\cite{wang2022transtab} trains transformer architectures to model semantic linkages among columns by taking column names into account, rather than relying solely on cell values. The aggregated knowledge about numerous tables and columns supports both zero-shot and few-shot inference when encountering novel tabular problems. Similarly, TabPFN~\cite{hollmann2022tabpfn} generates synthetic datasets by sampling diverse distributions, facilitating pre-training of Bayesian neural networks. After this pre-training phase, both the train and test examples from the target task are jointly inputted at inference time, obviating the need for further training. Insights transferred from broad data domains allow these methods to outperform classic tree-based algorithms and prove particularly advantageous under few-shot regimes.

\subsection{Language-Interfaced Tabular Learning}  
Leveraging the capabilities of large language models (LLMs)~\cite{anil2023palm,openai2023gpt4} represents another fruitful strategy for harnessing prior knowledge. By virtue of training on vast and heterogeneous text datasets, LLMs exhibit a strong capacity to generalize and adapt to new, previously unseen applications across different areas~\cite{brown2020language}. Recent efforts have exploited the extensive prior understanding embedded in LLMs to facilitate tabular learning. This typically involves serializing the tabular data as text and presenting it through input prompts~\cite{dinh2022lift,hegselmann2023tabllm,wang2023anypredict}. Prompts are designed to include representations of the data, features, and descriptions of the target task, enabling LLMs to interpret relationships among features and tasks and to make predictions. Research trends have shifted towards tuning LLMs for tabular tasks with efficient parameter optimization strategies such as LoRA~\cite{hu2021lora} or IA3~\cite{liu2022few}. Alternatively, in-context learning has been popular, wherein selected few-shot demonstrations are included directly in the prompts to elicit task-specific reasoning—often without any gradient-based updates~\cite{dinh2022lift,hegselmann2023tabllm,nam2023semi,slack2023tablet}.

Although the previously discussed techniques have demonstrated success in enhancing few-shot tabular learning, they rely on routing inference through the LLM, which can present significant obstacles in industrial settings. This paper proposes moving away from the standard black-box paradigm in which every individual prediction is made end-to-end via the LLM. Instead, the method focuses on using the LLM to explicitly derive the rules or decision criteria tied to each class label. \textsf{FeatLLM} translates these extracted rules into programmatic representations that are used to engineer additional features. This enables the prediction process to bypass the LLM by leveraging features constructed from rules identified through in-context learning on limited labeled data and pre-existing knowledge.

\section{Methods}
\textsf{FeatLLM} is intended to identify the "rules"—that is, the feature-based conditions distinguishing each possible output class—from a small set of provided instances, as illustrated in Figure~\ref{fig:main_model}. These derived rules, generated by the LLM, are translated into binary features which signify whether or not a given rule holds true for a specific sample. These binary indicators are then combined via a weighted dot product (with the weights constrained to be non-negative and learned during training), producing a class probability estimate for the sample. Through this supervised training procedure, the model acquires knowledge of which features are most predictive for each class (see Section~\ref{Sec:training} for details). The method is further enhanced by performing this entire rule extraction and prediction pipeline multiple times across different subsets of the data (i.e., via bagging), with ensemble averaging used to form the final prediction. The subsequent subsections detail the problem setup and each methodological component.

\vspace*{-0.12in}
\paragraph{Problem formulation.} Consider a tabular dataset comprising $N$ labeled instances $\mathcal{D} = \{(\mathbf{x}^i, \mathbf{y}^i)\}_{i=1}^{N}$. Each data point $\mathbf{x}^i$ is a $d$-dimensional feature vector, corresponding to $d$ distinct features, and is annotated with a class label $\mathbf{y}^i$ drawn from a fixed set of categories in a classification context. The associated feature names are specified as $F = \{f_j\}_{j=1}^{d}$, each accompanied by its natural language definition. In the $k$-shot learning scenario, the model is trained using only $k$ ($<N$) randomly chosen labeled examples from $\mathcal{D}$.

\subsection{Prompt Design for Extracting Rules}
\label{Sec:prompt_design}
In order to direct the LLM to derive rules that follow more reliable reasoning paths, we structure the prompt to reflect the logical process an expert would utilize when confronted with a tabular prediction task. As depicted in Fig.~\ref{fig:prompt_for_rule} and detailed in Appendix~\ref{sec:example_rule}, the prompt consists of the following three core elements:

\vspace*{-0.12in}
\paragraph{Basic information description.} This segment presents fundamental details necessary for addressing the task (highlighted in orange in Fig.~\ref{fig:prompt_for_rule}). It encompasses a clear articulation of the problem along with the class label information, definitions for the relevant features, and selected illustrative examples. The problem statement is cast as a direct inquiry, such as: ``Does this patient's myocardial infarction complications data indicate chronic heart failure? Yes or no?". Additional instances can be found in Appendix~\ref{sec:task_desc}. Feature descriptions clarify the data type for each attribute and may also provide definitions; for categorical attributes, a set of example category values can be listed. The demonstrations include a limited number of labeled training samples, transformed into textual representations consistent with serialization approaches in earlier works~\cite{dinh2022lift,hegselmann2023tabllm}:

\begin{align}
&\text{Serialize}(\mathbf{x}^i,\mathbf{y}^i, F) = \nonumber \\ 
&``\ f_1\ \text{is}\ \mathbf{x}^i_1.\ ... \ f_d\  \text{is}\  \mathbf{x}^i_d.\ \ \text{Answer:}\  \mathbf{y}^i\ ”
\end{align}

\vspace*{-0.12in}
\paragraph{Reasoning instruction.} Our objective is to improve the reasoning capabilities of LLMs by supplying them with explicit reasoning instructions (refer to blue highlights in Fig.~\ref{fig:prompt_for_rule}). This begins with an opening statement akin to the chain-of-thought paradigm~\cite{wei2022chain}, after which the inference task is decomposed into two stages. Initially, the LLM is prompted to hypothesize potential causal links or propensities between various features and the given task definition, entirely independent of the demonstration examples. This process leverages the model's general and commonsense knowledge, ensuring extraction of a comprehensive set of prior assumptions relevant to the task. In the subsequent stage, the LLM integrates example-based demonstrations alongside insights gathered from the first phase to establish class-specific rules. Segmenting reasoning in this fashion helps mitigate the formation of spurious associations involving unimportant attributes, thereby directing attention to the most pertinent features. To avoid issues stemming from an insufficient or excessive number of rules—which can result in underfitting or unwieldy prompt lengths—a fixed quota (namely, 10)\footnote{For an empirical investigation of rule quantity, see Section~\ref{sec:analysis}.} is enforced per class. Furthermore, when crafting these rules, the LLM is instructed to utilize information about each feature's data type from the fundamental description to preclude logical inconsistencies arising from type errors.

\vspace*{-0.12in}
\paragraph{Response instruction.} To streamline the interpretation and subsequent application of rules inferred by the LLM, explicit directives are given regarding response formatting (see yellow highlights in Fig.~\ref{fig:prompt_for_rule}). The prompt details both the necessary response template for each target class (i.e., answer format) and specifies the formal structure that individual rules should adopt (i.e., rule syntax). These instructions equip the LLM to align its rule format with the particular characteristics of the feature types involved. When no additional specification is given, the LLM has flexibility to aggregate rules utilizing logical conjunctions such as AND or OR. A concrete example of the expected response can be found in Appendix~\ref{sec:outcome-rule}.

\subsection{Inferring Class Likelihood via Rules}
\label{Sec:training}

\paragraph{Parsing rules for feature generation.} In this step, we employ the rules derived from the earlier phase to construct novel binary features. For every class, we produce a feature reflecting whether a specific sample complies with the corresponding class rules (yielding either '0' or '1'). These features serve to quantify the likelihood that a sample pertains to a given class. Nonetheless, as the rules formulated by the LLM are articulated in natural language, there is a need to establish a suitable parsing approach to automate the generation of these features. While response and rule formats are deliberately restricted during rule generation to enhance parseability, LLM outputs are susceptible to variations in syntax~\cite{kovavc2023large}. As an example, the model may replace operators such as ``$>$” and ``$<$” with textual equivalents like ``greater than" or ``less than," introducing further inconsistency and complicating parsing procedures.

To overcome the difficulties posed by noisy natural language parsing, we refrain from constructing intricate programmatic parsing mechanisms. Instead, we leverage the LLM itself to handle this variability. Owing to its capacity for semantic interpretation and proficiency in generating programmatic solutions (as it is trained on code as well as text), the LLM can effectively create relevant code when supplied with the function name, descriptions of inputs and outputs, along with the inferred rules within the prompt. The prompts and their resultant code outputs can be found in Figures~\ref{fig:prompt_for_parsing} and ~\ref{sec:example_parsing}-\ref{sec:example_parse_outcome} in the Appendix. This generated code is then run with Python’s \texttt{exec()} function for converting the data as needed.

\vspace*{-0.12in}
\paragraph{Inferring class likelihood.} Upon obtaining the set of rule-based binary features for each class, the simplest approach to estimate the likelihood that a sample belongs to a particular class is by tallying the number of satisfied rules for that class (i.e., summing the associated binary features). Nevertheless, given that individual rules or features may exhibit varying levels of importance, their relative weights should ideally be inferred using labeled training data. To capture these distinctions, we employ a standard machine learning (ML) algorithm that operates over each class’s binary feature vector. As an illustration, we can use a linear model excluding bias. Let $\mathbf{z}_k^i \in \{0, 1\}^R$ denote the binary feature vector derived from sample $\mathbf{x}^i$ with respect to class $k$, where $R$ denotes the number of rules for each class and $c$ represents the total number of classes. The probability vector for each class, $\mathbf{p}^i$, is then calculated by:

\begin{align}
\text{logit}_k^i &= \max(\mathbf{w}_k, 0) \cdot \mathbf{z}_k^i, \nonumber \\
\mathbf{p}^i &= \text{Softmax}({[\text{logit}_{1}^i, ..., \text{logit}_{c}^i]}). \label{eq:infer_prob}
\end{align}

In this context, $\mathbf{w}_k$ denotes the adjustable parameters within the linear model, each reflecting the relative importance attributed to a corresponding feature. By constraining these weights to lie within the positive domain, we guarantee that the features proposed by the LLM cannot negatively impact the probability of assigning a class during training. The resulting logit vector undergoes transformation into class probabilities via the application of the softmax function. The weight parameters $\mathbf{w}_k$ are optimized by minimizing the cross-entropy loss, utilizing a limited set of few-shot training instances.

\vspace*{-0.12in}
\paragraph{Ensembling with bagging.} To conclude, we iteratively repeat the entire process, thereby constructing several inference models whose outputs are combined in an ensemble to obtain the final decision. For $T$ distinct runs, let the learned weights be $\{\mathbf{w}[t]\}_{t=1}^T$; the final predicted probability for each class, $\mathbf{p}^i$, is then computed by averaging the outputs generated using each individual $\mathbf{w}[t]$ as described in Eq.~\ref{eq:infer_prob}.

Ensemble diversity is promoted by setting a sufficiently elevated temperature during the LLM inference phase, which increases the stochasticity in rule creation. Furthermore, we permute the order of few-shot exemplars, motivated by findings that their sequence can alter the LLM's responses~\cite{lu2022fantastically}\footnote{A discussion regarding the diversity of rules is presented in Appendix~\ref{sec:diversity}}. To further enrich ensemble diversity and improve robustness, bagging is employed to subsample either features or training instances for each iteration, an approach that becomes increasingly valuable as the size of the dataset or feature set grows. This ensemble framework yields two principal benefits. Firstly, errors from LLM-derived rules in individual runs may be counterbalanced by accurate rules obtained in others, thereby reinforcing the stability of the overall method. This phenomenon is related to the self-consistency property of LLMs~\cite{wang2022self}, wherein rules that emerge recurrently across several trials tend to be more trustworthy. Secondly, the ensembling mechanism provides a solution to the constraints of LLM input context size. In circumstances where the tabular input would exceed the model's prompt window, bagging limits prompt length, thus maintaining system robustness. While a single rule set may omit certain feature dependencies, aggregating multiple weak learners, each generated from different trials, allows the ensemble to effectively mitigate the risk of incomplete feature or instance representation in any single trial.

\section{Experiments}
We assess the effectiveness of \textsf{FeatLLM} across a range of tabular datasets and carry out a variety of analyses to explore the mechanisms underpinning our framework. For reasons of brevity, supplementary studies—including experiments with alternative LLMs, in-depth qualitative evaluations, and other aspects—are provided in the Appendix.

\subsection{Performance Evaluation}
\paragraph{Datasets.}
Our study involves 13 datasets covering binary and multi-class classification scenarios: (1) Adult~\cite{asuncion2007uci}, aimed at forecasting whether an individual’s income surpasses \$50,000 per year; (2) Bank~\cite{moro2014data}, to anticipate if a client will subscribe to a term deposit; (3) Blood~\cite{yeh2009knowledge}, used for predicting if a donor is likely to return; (4) Car~\cite{kadra2021well}, designed to rate the quality category of a vehicle; (5) Communities~\cite{misc_communities_and_crime_183}, to forecast the crime rate within a locality; (6) Credit-g~\cite{kadra2021well}, which estimates whether someone constitutes a good or bad credit risk; (7) Diabetes\footnote{\texttt{kaggle.com/datasets/uciml/pima-indians-diabetes-database}}, concerning the likelihood of diabetes diagnosis for an individual; (8) Heart\footnote{\texttt{kaggle.com/datasets/fedesoriano/heart-failure-prediction}}, which detects the potential presence of coronary artery disease; and (9) Myocardial~\cite{misc_myocardial_infarction_complications_579}, focused on determining the possibility of chronic heart failure in patients.

Additionally, we introduce several contemporary and synthetic datasets that are seldom examined and were not included in the LLMs’ pre-training: (10) Cultivars, which estimates whether a particular soybean cultivar will result in high grain yield\footnote{\texttt{archive.ics.uci.edu/dataset/913}}; (11) NHANES, to classify a participant’s age group based on given data\footnote{\texttt{archive.ics.uci.edu/dataset/887}}; (12) Sequence-type, which labels each sequence as arithmetic, geometric, Fibonacci, or Collatz; and (13) Solution-mix, where the objective is to ascertain if the concentration in a mixture of four solutions exceeds 0.5. The first two of these have only been made available on the UCI Machine Learning Repository after September 2023, confirming their exclusion from the pre-training data of the LLM API (gpt-3.5-turbo-0613) employed. The last two are artificial datasets, eliminating any risk that the LLM memorized them during pre-training.

The datasets selected represent a spectrum of sizes and complexities, with further specifics in Table~\ref{Tab:basic-desc}. Each dataset comprises explicit attribute names and accompanying descriptions. Particularly, datasets such as `Communities' and `Myocardial' possess over 100 attributes each, thereby presenting difficulties related to the constrained context window capacity of LLMs.

\vspace*{-0.12in}
\paragraph{Baselines.}
For our comparisons, we benchmark against nine baseline approaches. The initial set encompasses traditional models for tabular data: (1) Logistic Regression (LogReg), (2) XGBoost~\cite{chen2016xgboost}, and (3) RandomForest~\cite{ho1995random}. The subsequent two employ unlabeled data, namely (4) SCARF~\cite{bahri2022scarf} and (5) STUNT~\cite{nam2023stunt}. It's worth highlighting the practical challenges of acquiring large-scale unlabeled tabular datasets. A more recent baseline, (6) TabPFN~\cite{hollmann2022tabpfn}, leverages synthetically generated datasets to diversify model pre-training. The final trio involves LLM-based solutions: (7) In-context learning (In-context)~\cite{wei2022emergent}, (8) TABLET~\cite{slack2023tablet}, and (9) TabLLM~\cite{hegselmann2023tabllm}. In-context learning infuses a small set of training exemplars directly within the input prompt, requiring no model adaptation. TABLET supplements the LLM’s prompt with additional information derived from an external classifier utilizing rule-based and prototype methods to improve inferential quality. TabLLM utilizes the IA3~\cite{liu2022few} technique for parameter-efficient tuning, adapting the LLM to the few-shot samples for fine-tuned performance.

\vspace*{-0.12in}
\paragraph{Implementation details.}
\textsf{FeatLLM} is built upon GPT-3.5 as its primary LLM foundation, but its framework is intentionally designed to remain independent of any specific LLM (see Appendix~\ref{sec:backbones} for evaluation using alternate backbones). During LLM inference, the temperature parameter is assigned a value of 0.5, and top-p retains its default API value of 1. For extracting rules, we utilize 20 ensembles and extract a total of 10 rules. Additional analysis of hyper-parameter effects is provided in Figure~\ref{fig:hyperparameter2} and Appendix~\ref{sec:hyper-parameter}.

For the linear model, training is conducted with the Adam optimizer at a learning rate of 0.01 over 200 epochs. Optimal epochs are chosen through $k$-fold cross-validation. In reproducing baseline methods, GPT-3.5 serves as the LLM for in-context learning baselines, while T0~\cite{sanh2022multitask} is adopted in TabLLM, as in the original configuration. It is important to mention that TabLLM limits LLM use to those with open-source checkpoints, hence GPT-3.5 cannot be employed in its experiments. For standard machine learning models such as LogReg, XGBoost, and RandomForest, their hyper-parameters are tuned by conducting grid search in conjunction with $k$-fold cross-validation. The value of $k$ is either 2 or 4, ensuring each class is present in the training split. For further implementation specifics, see Appendix~\ref{sec:baselines}.

\vspace*{-0.12in}
\paragraph{Results.}
Tables~\ref{tab:main1} and \ref{tab:main2} display the comparative performance results across various datasets. Each experiment was conducted three times with different random splits of the training and test sets; reported metrics represent the mean and standard deviation of the resulting AUC values. \textsf{FeatLLM} frequently achieves either the highest or second-highest ranking among all evaluated baseline models. Remarkably, our approach, with only 4 labeled examples per class, achieves results competitive with traditional tabular learning methods that utilize 32 examples (see Fig.~\ref{fig:compares}). Although STUNT and TabPFN show pronounced gains on individual datasets, their aggregate advantage over conventional machine learning algorithms remains limited. Additionally, LLM-powered methods—such as in-context learning, TABLET, and TabLLM—encounter failures when applied to datasets with extensive feature sets. In contrast, our method, which leverages both bagging and ensemble strategies, demonstrates robust and satisfactory performance even on high-dimensional tabular data. Notably, our bagging and ensemble paradigm could be transferred to other LLM-based models; however, doing so would necessitate doubling the per-sample inference workload, greatly amplifying computational demands and rendering such adaptations largely infeasible in practice.

\subsection{Analysis \& Discussion}
\label{sec:analysis}
\paragraph{Ablation study.} To investigate the contributions of principal components, we conduct ablation experiments targeting: (1) \textsf{-Tuning}: removing the weight tuning phase from the linear model and instead aggregating binary feature values to form the logit; (2) \textsf{-Ensemble}: bypassing the ensemble phase; (3) \textsf{-Description}: excluding feature description utilization; and (4) \textsf{-Reasoning}: skipping the Step-1 component of the reasoning instruction during rule creation.

Table~\ref{tab:ablation} reports the average AUC change due to these ablations across all datasets. For each variant, the absence of a major component uniformly leads to diminished performance. Notably, the positive effect of weight tuning becomes more pronounced as the number of shots grows; this trend suggests that rule importance estimation can be performed more accurately with larger datasets, enhancing the impact of tuning. Conversely, feature description and explicit reasoning instructions exert a greater effect when data is limited, underscoring the significance of efficiently harnessing LLM prior knowledge for low-shot learning scenarios. Finally, our ensemble approach consistently yields performance improvements, regardless of setting.

\vspace*{-0.12in}
\paragraph{Training \& inference time comparison.} We present a comparative analysis of the training and inference times among various approaches. These evaluations utilize the Adult dataset, employing a training set that consists of 16 samples. A single A100 GPU serves as the standard computational resource for all methods, apart from TabLLM, which necessitates four A100 GPUs due to its model parallelism requirements. The timing results are reported in Table~\ref{tab:cost}. Importantly, \textsf{FeatLLM} achieves inference times that are on par with those of traditional techniques, distinguishing itself from other LLM-based approaches, including In-context learning, TABLET, and TabLLM, which require substantially greater inference durations. When examining training duration, \textsf{FeatLLM}'s performance is consistent with other LLM-based strategies, involving only 30 API calls. This minimal requirement does not generate significant costs and is amenable to parallel execution.

\vspace*{-0.12in}
\paragraph{Quality of features generated by \textsf{FeatLLM}.}
To evaluate the informativeness of the rule-derived features produced by \textsf{FeatLLM}, we performed a straightforward study focusing on their utility in real-world tasks. In particular, we computed the mean AUROC value between the generated features and the ground-truth labels, subsequently calculating the proportion of features whose AUROC exceeded 0.5 within each dataset. The outcomes of this analysis are detailed in Table~\ref{tab:quality}. The findings indicate that most of the produced rules are pertinent to the target variable, thus offering meaningful contributions to solving the main task (see the ``Before tuning'' column). Furthermore, as the linear model is fine-tuned using data, features of marginal utility—those for which the corresponding learned weights fall beneath a predefined threshold—are automatically pruned (see the ``After tuning'' column). The threshold applied here is 0.1, as determined based on the aggregate distribution of weight values.

\vspace*{-0.12in}
\paragraph{Effect of adding spurious correlations.} We performed an investigation to evaluate how different approaches handle spurious, noisy correlations when available labeled data is extremely limited. Specifically, the Heart dataset was augmented with randomly selected columns from the Adult dataset, which have no relevance to the Heart dataset’s objective. We assessed how the predictive performance of each model changed as we systematically increased the number of these unrelated features. For fairness, the experiments utilized a consistent sample of 8 labeled data points per task and excluded the use of any unlabeled data, thereby omitting algorithms such as SCARF and STUNT from this benchmark.

In Figure~\ref{fig:spurious}, the effect of introducing irrelevant features on model accuracy is depicted. Notably, \textsf{FeatLLM} experienced the smallest decline in accuracy as spurious correlations increased, outperforming comparison methods. This resilience is plausibly attributed to our method’s reliance on prior knowledge about feature-task relations during rule extraction, which discourages forming rules that depend on non-informative columns and enhances robustness. In practice, rules involving noisy features constituted merely 1–2\% of the overall rules. However, it is also evident that \textsf{FeatLLM} may be susceptible to incorporating spurious associations or misattributing causality if randomly injected columns are derived from datasets closely aligned to the task (refer to Appendix~\ref{sec:spurious-appendix}).

\vspace*{-0.12in}
\paragraph{Hyper-parameter analysis}
\textsf{FeatLLM} employs two primary hyper-parameters: the ensemble size and the number of rules per class retrieved in each LLM invocation. To assess how these choices affect model effectiveness, we ran a series of experiments. Our observations indicate that increasing the ensemble size correlates with improved performance; however, beyond a threshold (approximately 20 ensembles), these gains plateau (refer to Fig.~\ref{fig:appendix-hyperparameter1} in the Appendix). Regarding the number of rules, although variations did not yield statistically significant changes in results, we found that setting this value to 10 optimally balances prompt length against predictive performance (see Fig.~\ref{fig:hyperparameter2}). We hypothesize that while more rules augment the information by expanding input dimensionality, this may cause overfitting when available data is sparse.

\section{Conclusion}
In this work, we introduce \textsf{FeatLLM}, a framework that establishes a new benchmark for few-shot tabular learning utilizing LLMs. Rather than directly producing per-sample predictions, \textsf{FeatLLM} leverages LLMs as automated feature engineers, deriving decision criteria for predictions. Through a rule extraction mechanism and an ensemble strategy using bagging, \textsf{FeatLLM} achieves competitive results with minimal inference overhead, requires only access to an LLM API (no model retraining), and circumvents challenges associated with large feature spaces. Notably, \textsf{FeatLLM} combines efficiency with predictive strength, making it a robust solution for real-world tabular problems with limited annotated data.

The current scope of \textsf{FeatLLM} addresses the low-shot learning scenario. Moving forward, we aim to broaden the framework to generate novel features for larger datasets, explore a wider variety of feature forms beyond decision rules, and incorporate interpretability, thus facilitating its utility across more diverse problem domains.

\section{Broader Impact}
\textsf{FeatLLM} is designed to tap into the extensive prior knowledge encapsulated within LLMs, enabling advances in tabular learning beyond what traditional supervised algorithms can offer. By improving data efficiency, our research addresses key overfitting challenges in situations with limited training samples, using LLMs to inject external information. This framework is especially valuable for practitioners in practical domains like finance or healthcare, where data annotation is either costly or logistically difficult, and where LLMs pretrained on relevant corpora provide critical supplemental knowledge. As \textsf{FeatLLM} gains wider traction, a significant future consideration will be systematically studying how inherent societal biases or misinformation in LLM prior knowledge might influence predictions, potentially overriding observed labeled examples and necessitating careful mitigation strategies.

\end{document}